% Generated by docs/guidebook/gen_tables.py from the code. Do not edit by hand.
\section{Build the spec}
\begin{toolentry}{new\_spec}{name, size, length\_unit, force\_unit}{profiles: spec, clarify, full; changes the spec (undoable)}\label{tool:new_spec}\index{new_spec@\texttt{new\_spec} (tool)}
Start a new spec: a box design domain of the given size (x, y, z) with its min corner at the origin. Replaces any current spec. Material defaults to steel (E 210e3, nu 0.3) until set\_params changes it.
\par\smallskip\begin{argtable}
\code{name} & string & \emph{required} \\
\code{size} & [x, y, z] & \emph{required} \\
\code{length\_\allowbreak{}unit} & \code{mm} | \code{m} | \code{in} & \code{"mm"} \\
\code{force\_\allowbreak{}unit} & \code{N} | \code{kN} | \code{lbf} & \code{"N"} \\
\end{argtable}
\end{toolentry}
\begin{toolentry}{set\_spec}{spec}{profiles: spec, clarify, full; changes the spec (undoable)}\label{tool:set_spec}\index{set_spec@\texttt{set\_spec} (tool)}
Replace the whole spec with an agentto.spec/v0 document. A document that does not pass the schema yet (e.g. no load) is kept as a draft and its schema errors are returned.
\par\smallskip\begin{argtable}
\code{spec} & object & \emph{required} \\
\end{argtable}
\end{toolentry}
\begin{toolentry}{load\_template}{name, k}{profiles: full; changes the spec (undoable)}\label{tool:load_template}\index{load_template@\texttt{load\_template} (tool)}
Load a FEAD benchmark problem (2:1:1 box, E = 1) at k elements on the short edge.
\par\smallskip\begin{argtable}
\code{name} & \code{cantilever} | \code{bridge} | \code{torsion} & \emph{required} \\
\code{k} & integer & \code{40} \\
\end{argtable}
\end{toolentry}
\begin{toolentry}{set\_params}{name, description, size, length\_unit, force\_unit, E, nu, material\_name, volume\_fraction, h, cells\_short\_edge, min\_feature, clear\_min\_feature}{profiles: spec, clarify, full; changes the spec (undoable)}\label{tool:set_params}\index{set_params@\texttt{set\_params} (tool)}
Set scalar spec fields; only the arguments given change. E is in force/length\textasciicircum{}2 of the spec units (210 GPa = 210e3 N/mm\textasciicircum{}2). volume\_fraction is of the design (non-passive) volume, 0-1. Resolution: give h (element edge, length units) or cells\_short\_edge, not both.
\par\smallskip\begin{argtable}
\code{name} & string, optional & -- \\
\code{description} & string, optional & -- \\
\code{size} & [x, y, z], optional & -- \\
\code{length\_\allowbreak{}unit} & \code{mm} | \code{m} | \code{in}, optional & -- \\
\code{force\_\allowbreak{}unit} & \code{N} | \code{kN} | \code{lbf}, optional & -- \\
\code{E} & number, optional & -- \\
\code{nu} & number, optional & -- \\
\code{material\_\allowbreak{}name} & string, optional & -- \\
\code{volume\_\allowbreak{}fraction} & number, optional & -- \\
\code{h} & number, optional & -- \\
\code{cells\_\allowbreak{}short\_\allowbreak{}edge} & integer, optional & -- \\
\code{min\_\allowbreak{}feature} & number, optional & -- \\
\code{clear\_\allowbreak{}min\_\allowbreak{}feature} & boolean & \code{false} \\
\end{argtable}
\end{toolentry}
\begin{toolentry}{add\_support}{id, region, fix, note}{profiles: spec, clarify, full; changes the spec (undoable)}\label{tool:add_support}\index{add_support@\texttt{add\_support} (tool)}
Fix the nodes of a region in the given directions. Regions: face \{face\}; patch \{face, center, size, shape rect\textbar{}ellipse\} with center/size in the face's two in-plane axes in increasing order ((y,z) on x faces, (x,z) on y faces, (x,y) on z faces); edge \{axis, at, span?\}; box \{min, max\}; cylinder \{axis, center, radius, span\}. Coordinates are absolute, origin at the domain's min corner.
\par\smallskip\begin{argtable}
\code{id} & string & \emph{required} \\
\code{region} & region (face, patch, edge, box, cylinder) & \emph{required} \\
\code{fix} & list of \code{x} | \code{y} | \code{z} & \code{["x",\allowbreak{} "y",\allowbreak{} "z"]} \\
\code{note} & string & \code{""} \\
\end{argtable}
\end{toolentry}
\begin{toolentry}{add\_load}{id, region, force, note}{profiles: spec, clarify, full; changes the spec (undoable)}\label{tool:add_load}\index{add_load@\texttt{add\_load} (tool)}
Apply a total force vector (spec force units, x y z components) shared equally by the region's nodes. Region forms as in add\_support.
\par\smallskip\begin{argtable}
\code{id} & string & \emph{required} \\
\code{region} & region (face, patch, edge, box, cylinder) & \emph{required} \\
\code{force} & [x, y, z] & \emph{required} \\
\code{note} & string & \code{""} \\
\end{argtable}
\end{toolentry}
\begin{toolentry}{add\_passive}{id, kind, region, note}{profiles: spec, clarify, full; changes the spec (undoable)}\label{tool:add_passive}\index{add_passive@\texttt{add\_passive} (tool)}
Keep-in (solid) or keep-out (void) volume, a box or cylinder: elements whose centres fall inside are held at that state and excluded from the volume fraction.
\par\smallskip\begin{argtable}
\code{id} & string & \emph{required} \\
\code{kind} & \code{solid} | \code{void} & \emph{required} \\
\code{region} & region (box, cylinder) & \emph{required} \\
\code{note} & string & \code{""} \\
\end{argtable}
\end{toolentry}
\begin{toolentry}{update\_item}{id, region, fix, force, kind, note, new\_id}{profiles: spec, clarify, full; changes the spec (undoable)}\label{tool:update_item}\index{update_item@\texttt{update\_item} (tool)}
Change fields of an existing support, load or passive region (only the arguments given).
\par\smallskip\begin{argtable}
\code{id} & string & \emph{required} \\
\code{region} & object, optional & -- \\
\code{fix} & list of \code{x} | \code{y} | \code{z}, optional & -- \\
\code{force} & [x, y, z], optional & -- \\
\code{kind} & \code{solid} | \code{void}, optional & -- \\
\code{note} & string, optional & -- \\
\code{new\_\allowbreak{}id} & string, optional & -- \\
\end{argtable}
\end{toolentry}
\begin{toolentry}{remove\_item}{id}{profiles: spec, clarify, full; changes the spec (undoable)}\label{tool:remove_item}\index{remove_item@\texttt{remove\_item} (tool)}
Delete a support, load or passive region by id.
\par\smallskip\begin{argtable}
\code{id} & string & \emph{required} \\
\end{argtable}
\end{toolentry}
\begin{toolentry}{undo}{}{profiles: spec, clarify, full; changes the spec (undoable)}\label{tool:undo}\index{undo@\texttt{undo} (tool)}
Revert the last spec change.
\par\smallskip{\small\sffamily No arguments.}
\end{toolentry}
\begin{toolentry}{redo}{}{profiles: spec, clarify, full; changes the spec (undoable)}\label{tool:redo}\index{redo@\texttt{redo} (tool)}
Re-apply the last spec change that undo reverted.
\par\smallskip{\small\sffamily No arguments.}
\end{toolentry}
\section{Inspect and check}
\begin{toolentry}{get\_spec}{}{profiles: spec, clarify, full}\label{tool:get_spec}\index{get_spec@\texttt{get\_spec} (tool)}
The current spec document (possibly incomplete).
\par\smallskip{\small\sffamily No arguments.}
\end{toolentry}
\begin{toolentry}{list\_regions}{h}{profiles: spec, clarify, full}\label{tool:list_regions}\index{list_regions@\texttt{list\_regions} (tool)}
Ground the spec on the mesh: domain faces and their in-plane axes, then every support/load/passive item with its physical bounding box and size, and how many nodes (or elements) it selects at mesh size h (default: the spec's resolution), with any snapping. An item that selects 0 nodes does nothing.
\par\smallskip\begin{argtable}
\code{h} & number, optional & -- \\
\end{argtable}
\end{toolentry}
\begin{toolentry}{validate\_spec}{h}{profiles: spec, clarify, full}\label{tool:validate_spec}\index{validate_spec@\texttt{validate\_spec} (tool)}
Run the schema and the deterministic rule base (well-posedness, rigid-body modes, empty regions, loads on fixed nodes, overlaps, resolvable minimum feature, memory) at mesh size h (default: the spec's). Returns ok plus findings with stable codes and suggested fixes. Only an ok spec can be solved.
\par\smallskip\begin{argtable}
\code{h} & number, optional & -- \\
\end{argtable}
\end{toolentry}
\begin{toolentry}{explain\_spec}{}{profiles: spec, clarify, full}\label{tool:explain_spec}\index{explain_spec@\texttt{explain\_spec} (tool)}
Plain-language summary of the spec, for confirming intent with the user.
\par\smallskip{\small\sffamily No arguments.}
\end{toolentry}
\begin{toolentry}{render}{view}{profiles: spec, clarify, full}\label{tool:render}\index{render@\texttt{render} (tool)}
PNG picture of the domain with supports (blue), load arrows (red) and keep-in/keep-out regions, drawn from the spec geometry. Use it to confirm the setup with the user before solving.
\par\smallskip\begin{argtable}
\code{view} & \code{iso} | \code{front} | \code{top} | \code{side} | \code{back\_\allowbreak{}iso} & \code{"iso"} \\
\end{argtable}
\end{toolentry}
\begin{toolentry}{ladder}{}{profiles: spec, clarify, full}\label{tool:ladder}\index{ladder@\texttt{ladder} (tool)}
Resolution rungs for this spec (sketch, preview, refine, release; each 2x finer per axis). The spec's resolution h is the base mesh the ladder is built from, not the release mesh: release is the finest rung that fits on the GPU, often finer than the spec's h. Only a solve at the release rung can be released.
\par\smallskip{\small\sffamily No arguments.}
\end{toolentry}
\section{Ask the user}
\begin{toolentry}{ask\_clarification}{question, options, previews, about}{profiles: clarify, full}\label{tool:ask_clarification}\index{ask_clarification@\texttt{ask\_clarification} (tool)}
Ask the user ONE question when the brief is ambiguous in a way that changes the design (which face, how large, what force or material). Give 2-4 options, each a concrete reading: the full item the spec would contain (item\_kind + item, as you would pass to add\_support / add\_load / add\_passive) and/or parameter values (params, e.g. \{"volume\_fraction": 0.3\}). Returns the chosen option, which you then apply with the editing tools, or choice null with the user's comment. previews=true (full profile) first solves each reading at the sketch rung, so the user sees what it does to the design. about: the spec items the question concerns (keys as for set\_source); in a provenance session every one must be tagged assumed.
\par\smallskip\begin{argtable}
\code{question} & string & \emph{required} \\
\code{options} & list of options (see below) & \emph{required} \\
\code{previews} & boolean & \code{false} \\
\code{about} & list of string, optional & -- \\
\end{argtable}
\end{toolentry}
\paragraph{Clarification options.} Each element of \code{options} is an object with these fields (\code{ClarifyOption}); unknown fields are rejected.
\begin{argtable}
\code{label} & string & \emph{required} \\
\code{item\_\allowbreak{}kind} & \code{support} | \code{load} | \code{passive}, optional & -- \\
\code{item} & object, optional & -- \\
\code{params} & object, optional & -- \\
\end{argtable}
\section{Record value sources}
\begin{toolentry}{set\_source}{key, source, note}{profiles: full}\label{tool:set_source}\index{set_source@\texttt{set\_source} (tool)}
Record where one spec item's value comes from. key: units, domain, material, volume\_fraction, resolution, or support:<id>, load:<id>, passive:<id>. source: stated (the brief gives it), derived (computed from numbers the brief gives), assumed (your choice where the brief is silent or ambiguous), answered (fixed by an answer to ask\_clarification). note: what was stated, derived or assumed.
\par\smallskip\begin{argtable}
\code{key} & string & \emph{required} \\
\code{source} & \code{stated} | \code{derived} | \code{assumed} | \code{answered} & \emph{required} \\
\code{note} & string & \code{""} \\
\end{argtable}
\end{toolentry}
\begin{toolentry}{list\_sources}{}{profiles: full}\label{tool:list_sources}\index{list_sources@\texttt{list\_sources} (tool)}
Every spec item with its source tag (stated, derived, assumed, answered), the untagged items, and the items tagged assumed.
\par\smallskip{\small\sffamily No arguments.}
\end{toolentry}
\section{Solve}
\begin{toolentry}{solve}{rung, h, warm, optimize, wait, timeout\_s}{profiles: full}\label{tool:solve}\index{solve@\texttt{solve} (tool)}
Optimize the current spec on a ladder rung (or mesh size h). Warm-starts from this spec's design on the 2x coarser grid when there is one. By default that warm start only certifies the upsampled coarse design (about 12 iterations); optimize=True re-opens it (about 60 iterations, \textasciitilde{}60\% of a cold run, compliance within +0.1\% of cold or better on the X1 problems). wait=False returns the job id at once. The result carries the solver's certificate; only a certified release is a released design.
\par\smallskip\begin{argtable}
\code{rung} & \code{sketch} | \code{preview} | \code{refine} | \code{release}, optional & -- \\
\code{h} & number, optional & -- \\
\code{warm} & boolean & \code{true} \\
\code{optimize} & boolean & \code{false} \\
\code{wait} & boolean & \code{true} \\
\code{timeout\_\allowbreak{}s} & number & \code{900} \\
\end{argtable}
\end{toolentry}
\begin{toolentry}{status}{}{profiles: full}\label{tool:status}\index{status@\texttt{status} (tool)}
Jobs of this session, newest last: rung h, status, iteration, certified.
\par\smallskip{\small\sffamily No arguments.}
\end{toolentry}
\begin{toolentry}{wait\_job}{job, timeout\_s}{profiles: full}\label{tool:wait_job}\index{wait_job@\texttt{wait\_job} (tool)}
Wait for a queued or running solve (default: the newest job) and return its result, as solve(wait=True) would. Use this after a solve timed out instead of cancelling and resubmitting.
\par\smallskip\begin{argtable}
\code{job} & string, optional & -- \\
\code{timeout\_\allowbreak{}s} & number & \code{900} \\
\end{argtable}
\end{toolentry}
\begin{toolentry}{cancel}{}{profiles: full}\label{tool:cancel}\index{cancel@\texttt{cancel} (tool)}
Stop the running solve (it ends with status cancelled at the next iteration).
\par\smallskip{\small\sffamily No arguments.}
\end{toolentry}
\begin{toolentry}{field\_stats}{job, region, box\_min, box\_max}{profiles: full}\label{tool:field_stats}\index{field_stats@\texttt{field\_stats} (tool)}
Numbers on a finished design (default: the latest finished job): material fraction, grey fraction, solid components, and whether each load is connected through solid to a support. With region (an item id) or box\_min/box\_max, also the density inside that region.
\par\smallskip\begin{argtable}
\code{job} & string, optional & -- \\
\code{region} & string, optional & -- \\
\code{box\_\allowbreak{}min} & [x, y, z], optional & -- \\
\code{box\_\allowbreak{}max} & [x, y, z], optional & -- \\
\end{argtable}
\end{toolentry}
\section{Export}
\begin{toolentry}{export\_stl}{job, pool}{profiles: full}\label{tool:export_stl}\index{export_stl@\texttt{export\_stl} (tool)}
Write the design's rho = 0.5 surface as a binary STL (default: the latest finished job; pool > 1 averages pool\textasciicircum{}3 element blocks first). Only a certified release is a released design: anything else is exported as a draft and says so.
\par\smallskip\begin{argtable}
\code{job} & string, optional & -- \\
\code{pool} & integer & \code{1} \\
\end{argtable}
\end{toolentry}
\begin{toolentry}{export\_report}{job}{profiles: full}\label{tool:export_report}\index{export_report@\texttt{export\_report} (tool)}
Write the design report (JSON and Markdown): spec and hash, rung and solver, iterations and time, compliance, every certificate check, the STL if one was exported, and the hash of this session's log.
\par\smallskip\begin{argtable}
\code{job} & string, optional & -- \\
\end{argtable}
\end{toolentry}
